\chapter{Weighting Estimators}
\label{ch:weighting}

% ============================================================================
% Chapter 5: Weighting Estimators (IPW, DR, TMLE)
% Written: 2025-12-31
% Target: ~30 pages, ~850 lines
% ============================================================================

\section{Introduction}
\label{sec:weighting_intro}

Propensity score matching (Chapter~\ref{ch:propensity}) creates comparable groups by discarding unmatched units. \textbf{Weighting estimators} offer an alternative approach: rather than discarding units, we reweight them to create a pseudo-population where treatment is independent of covariates.

The fundamental insight is that inverse probability weighting transforms observational data into a pseudo-randomized experiment. If unit $i$ has treatment probability $e(X_i) = 0.2$, they represent 5 similar units who did not receive treatment---hence weight $1/0.2 = 5$.

\subsection{Chapter Overview}

This chapter covers three progressively more sophisticated weighting approaches:

\begin{enumerate}
    \item \textbf{Inverse Probability Weighting (IPW)}: Foundational weighting using only propensity
    \item \textbf{Doubly Robust (DR/AIPW)}: Combines IPW with outcome regression; consistent if \emph{either} model is correct
    \item \textbf{Targeted Maximum Likelihood (TMLE)}: Achieves semiparametric efficiency through iterative targeting
\end{enumerate}

\subsection{Assumptions}

All weighting estimators in this chapter require the same identification assumptions as propensity score matching:

\begin{assumption}[Unconfoundedness]
\begin{equation}
\poty{0}, \poty{1} \indep T \mid X
\end{equation}
\end{assumption}

\begin{assumption}[Positivity / Overlap]
\begin{equation}
0 < e(X) = \Prob(T = 1 \mid X) < 1 \quad \text{for all } X
\end{equation}
\end{assumption}

\begin{insight}
Positivity is \emph{critical} for weighting methods. When $e(X) \approx 0$ or $e(X) \approx 1$, weights become extreme and variance explodes. This chapter emphasizes practical strategies for diagnosing and addressing positivity violations.
\end{insight}

% ============================================================================
\section{Inverse Probability Weighting}
\label{sec:ipw}
% ============================================================================

\subsection{The IPW Estimator}

Under unconfoundedness, we cannot directly compare treated and control outcomes because treated units differ systematically from controls. IPW corrects this by weighting each unit inversely by their probability of receiving the treatment they actually received.

\begin{definition}[IPW Estimator for ATE]
The Horvitz-Thompson IPW estimator is:
\begin{equation}
\hat{\tau}_{\text{IPW}} = \frac{1}{n}\sum_{i=1}^{n} \left[ \frac{T_i Y_i}{e(X_i)} - \frac{(1-T_i) Y_i}{1 - e(X_i)} \right]
\label{eq:ipw_ate}
\end{equation}
\end{definition}

\begin{theorem}[Unbiasedness of IPW]
\label{thm:ipw_unbiased}
Under unconfoundedness and positivity, $\E[\hat{\tau}_{\text{IPW}}] = \tau_{\text{ATE}}$.
\end{theorem}

\begin{proof}
For the treated term:
\begin{align}
\E\left[\frac{T Y}{e(X)}\right] &= \E\left[\E\left[\frac{T Y}{e(X)} \mid X\right]\right] \\
&= \E\left[\frac{1}{e(X)} \E[T Y \mid X]\right] \\
&= \E\left[\frac{1}{e(X)} \cdot e(X) \cdot \E[Y \mid T=1, X]\right] \\
&= \E[\E[\poty{1} \mid X]] = \E[\poty{1}]
\end{align}
The last equality uses unconfoundedness: $\E[Y \mid T=1, X] = \E[\poty{1} \mid X]$.
Similarly, $\E[(1-T)Y/(1-e(X))] = \E[\poty{0}]$.
\end{proof}

\subsection{Stabilized Weights}

Standard IPW weights can have high variance when propensity scores are extreme. \textbf{Stabilized weights} multiply the standard weights by the marginal treatment probability:

\begin{definition}[Stabilized IPW Weights]
\begin{align}
w_i^{\text{stab}} &= \begin{cases}
\displaystyle\frac{\Prob(T=1)}{e(X_i)} & \text{if } T_i = 1 \\[2ex]
\displaystyle\frac{1 - \Prob(T=1)}{1 - e(X_i)} & \text{if } T_i = 0
\end{cases}
\end{align}
\end{definition}

Stabilized weights have mean approximately 1 (instead of potentially large values) while maintaining consistency.

\subsection{Weight Diagnostics}

Before trusting IPW estimates, examine the weight distribution:

\begin{definition}[Weight Diagnostics]
\begin{enumerate}
    \item \textbf{Weight range}: Check $[\min w_i, \max w_i]$; extreme values signal problems
    \item \textbf{Effective sample size}: Should be close to $n$:
    \[n_{\text{eff}} = (\textstyle\sum_i w_i)^2 / \sum_i w_i^2\]
    \item \textbf{Weight sum}: For balance, $\sum_{T=1} w_i \approx \sum_{T=0} w_i$
\end{enumerate}
\end{definition}

\begin{insight}
Rule of thumb: If any weight exceeds 20 (or maximum/minimum ratio exceeds 100), consider trimming or alternative estimators. Extreme weights dominate the estimate and inflate variance.
\end{insight}

\subsection{Propensity Trimming}

When positivity is violated, we can \emph{trim} units with extreme propensities:

\begin{equation}
\text{Keep unit } i \text{ if } \epsilon \leq \hat{e}(X_i) \leq 1 - \epsilon
\end{equation}

Common choices: $\epsilon \in \{0.01, 0.05, 0.10\}$.

\textbf{Trade-off}: Trimming reduces variance but introduces bias (changes the estimand from ATE to a conditional ATE for the trimmed population).

% ============================================================================
\section{Outcome Regression}
\label{sec:outcome_regression}
% ============================================================================

Before introducing doubly robust methods, we briefly review the pure outcome regression approach.

\begin{definition}[Outcome Regression Estimator]
Fit separate outcome models for treated and control:
\begin{align}
\hat{\mu}_1(X) &= \hat{\E}[Y \mid T=1, X] \\
\hat{\mu}_0(X) &= \hat{\E}[Y \mid T=0, X]
\end{align}
The ATE estimate is:
\begin{equation}
\hat{\tau}_{\text{OR}} = \frac{1}{n}\sum_{i=1}^{n} \left[ \hat{\mu}_1(X_i) - \hat{\mu}_0(X_i) \right]
\end{equation}
\end{definition}

Outcome regression is consistent if the outcome models $\mu_1, \mu_0$ are correctly specified. However, it can be sensitive to model misspecification, especially extrapolation in regions of poor overlap.

\subsection{Limitations}

\begin{itemize}
    \item \textbf{Model dependence}: Results depend on functional form assumptions
    \item \textbf{Extrapolation}: Predictions may be made far from observed data
    \item \textbf{No robustness}: If outcome model wrong, estimate is biased
\end{itemize}

These limitations motivate combining IPW with outcome regression.

% ============================================================================
\section{Doubly Robust Estimation}
\label{sec:doubly_robust}
% ============================================================================

\subsection{The AIPW Estimator}

The \textbf{Augmented Inverse Probability Weighted (AIPW)} estimator, also called the \textbf{Doubly Robust (DR)} estimator, combines IPW and outcome regression:

\begin{definition}[Doubly Robust Estimator]
\label{def:dr}
\begin{equation}
\hat{\tau}_{\text{DR}} = \frac{1}{n}\sum_{i=1}^{n} \left[ \hat{\psi}_1(O_i) - \hat{\psi}_0(O_i) \right]
\end{equation}
where:
\begin{align}
\hat{\psi}_1(O_i) &= \frac{T_i (Y_i - \hat{\mu}_1(X_i))}{\hat{e}(X_i)} + \hat{\mu}_1(X_i) \\[1ex]
\hat{\psi}_0(O_i) &= \frac{(1-T_i) (Y_i - \hat{\mu}_0(X_i))}{1 - \hat{e}(X_i)} + \hat{\mu}_0(X_i)
\end{align}
\end{definition}

The estimator has three components:
\begin{enumerate}
    \item $\hat{\mu}_1(X_i) - \hat{\mu}_0(X_i)$: outcome regression estimate
    \item $\frac{T_i(Y_i - \hat{\mu}_1(X_i))}{\hat{e}(X_i)}$: IPW correction for treated
    \item $\frac{(1-T_i)(Y_i - \hat{\mu}_0(X_i))}{1 - \hat{e}(X_i)}$: IPW correction for control
\end{enumerate}

\subsection{The Double Robustness Property}

\begin{theorem}[Double Robustness]
\label{thm:double_robust}
The DR estimator is consistent if \emph{either}:
\begin{enumerate}
    \item The propensity model $\hat{e}(X)$ is correctly specified, \textbf{or}
    \item The outcome models $\hat{\mu}_0(X), \hat{\mu}_1(X)$ are correctly specified.
\end{enumerate}
\end{theorem}

\begin{proof}[Proof sketch]
Consider the bias of the treated component. Let $e^*(X)$ denote the true propensity and $\mu_1^*(X)$ the true outcome model.

\textbf{Case 1}: Propensity correct ($\hat{e} = e^*$), outcome wrong.
The IPW term $\E[T(Y - \hat{\mu}_1)/e^*]$ equals $\E[\poty{1}] - \E[\hat{\mu}_1]$, and adding $\E[\hat{\mu}_1]$ gives $\E[\poty{1}]$. The outcome model errors cancel.

\textbf{Case 2}: Propensity wrong, outcome correct ($\hat{\mu}_1 = \mu_1^*$).
The residual $Y - \mu_1^*(X)$ has conditional mean zero given $X$ for treated units. The IPW term has expectation zero regardless of propensity specification. The outcome regression term provides the correct answer.

\textbf{Both correct}: DR estimator achieves the semiparametric efficiency bound (lowest possible variance among regular estimators).
\end{proof}

\begin{insight}
Double robustness provides an ``insurance policy'': you get two chances to correctly specify your models. If you're uncertain about functional forms, DR offers protection that neither IPW nor outcome regression alone provides.
\end{insight}

\subsection{When Both Models Are Wrong}

\textbf{Warning}: Double robustness does \emph{not} mean the estimator is robust to having \emph{both} models wrong. If both propensity and outcome models are misspecified, the DR estimator can be biased---sometimes more biased than IPW or outcome regression alone.

% ============================================================================
\section{Targeted Maximum Likelihood Estimation}
\label{sec:tmle}
% ============================================================================

TMLE improves on standard DR estimation by iteratively updating the outcome model to satisfy the efficient score equation, achieving the semiparametric efficiency bound with better finite-sample properties.

\subsection{Motivation}

Standard DR is a ``one-shot'' estimator: fit models, plug in, compute. TMLE adds a \textbf{targeting step} that updates the outcome model to:
\begin{enumerate}
    \item Reduce finite-sample bias
    \item Achieve the semiparametric efficiency bound
    \item Ensure valid influence function-based inference
\end{enumerate}

\subsection{The TMLE Algorithm}

\begin{definition}[TMLE for ATE]
\label{def:tmle_algorithm}
\textbf{Step 1: Initial Estimation}
\begin{itemize}
    \item Estimate propensity: $\hat{g}(X) = \Prob(T=1 \mid X)$
    \item Estimate outcome: $\hat{Q}(T, X) = \E[Y \mid T, X]$
\end{itemize}

\textbf{Step 2: Targeting}
\begin{itemize}
    \item Compute the \emph{clever covariate}:
    \begin{equation}
    H_i = \frac{T_i}{\hat{g}(X_i)} - \frac{1 - T_i}{1 - \hat{g}(X_i)}
    \label{eq:clever_covariate}
    \end{equation}
    \item Fit fluctuation model: regress $Y$ on $H$ with offset $\hat{Q}$
    \begin{equation}
    Y_i = \hat{Q}(T_i, X_i) + \varepsilon H_i + \text{error}
    \end{equation}
    \item Update predictions: $\hat{Q}^*(T, X) = \hat{Q}(T, X) + \hat{\varepsilon} H$
\end{itemize}

\textbf{Step 3: Iterate} until convergence (score equation satisfied):
\begin{equation}
\frac{1}{n}\sum_{i=1}^{n} H_i (Y_i - \hat{Q}^*(T_i, X_i)) \approx 0
\end{equation}

\textbf{Step 4: Estimate ATE}
\begin{equation}
\hat{\tau}_{\text{TMLE}} = \frac{1}{n}\sum_{i=1}^{n} \left[ \hat{Q}^*(1, X_i) - \hat{Q}^*(0, X_i) \right]
\end{equation}
\end{definition}

\subsection{The Clever Covariate}

The clever covariate $H$ plays a central role in TMLE. It is derived from the efficient influence function for the ATE:

\begin{equation}
\text{EIF}(O) = \frac{T(Y - \mu_1(X))}{e(X)} - \frac{(1-T)(Y - \mu_0(X))}{1 - e(X)} + \mu_1(X) - \mu_0(X) - \tau
\end{equation}

The targeting step ensures that $\E[\text{EIF}] = 0$, which guarantees efficient estimation.

\subsection{TMLE vs Standard DR}

\begin{center}
\begin{tabular}{lcc}
\toprule
Property & DR/AIPW & TMLE \\
\midrule
Double robustness & Yes & Yes \\
Semiparametric efficient & Sometimes & Always \\
Finite-sample bias & May be larger & Typically smaller \\
Computation & One-shot & Iterative \\
Bounded predictions & No & Can enforce \\
\bottomrule
\end{tabular}
\end{center}

\begin{insight}
TMLE is preferred when: (1) you need efficiency, (2) outcomes are bounded (e.g., probabilities), or (3) finite-sample properties matter. Standard DR is simpler and often suffices for large samples with good overlap.
\end{insight}

% ============================================================================
\section{Variance Estimation}
\label{sec:weighting_variance}
% ============================================================================

\subsection{Influence Function Approach}

All three estimators (IPW, DR, TMLE) can use influence function-based variance estimation:

\begin{definition}[Influence Function Variance]
For estimator $\hat{\tau}$ with influence function $\text{IF}_i$:
\begin{equation}
\widehat{\Var}(\hat{\tau}) = \frac{1}{n^2} \sum_{i=1}^{n} \text{IF}_i^2 = \frac{\widehat{\Var}(\text{IF})}{n}
\end{equation}
\end{definition}

\subsection{Influence Functions by Estimator}

\textbf{IPW Influence Function}:
\begin{equation}
\text{IF}_i^{\text{IPW}} = \frac{T_i Y_i}{\hat{e}(X_i)} - \frac{(1-T_i) Y_i}{1 - \hat{e}(X_i)} - \hat{\tau}_{\text{IPW}}
\end{equation}

\textbf{DR Influence Function}:
\begin{equation}
\text{IF}_i^{\text{DR}} = \hat{\psi}_1(O_i) - \hat{\psi}_0(O_i) - \hat{\tau}_{\text{DR}}
\end{equation}
where $\hat{\psi}_1, \hat{\psi}_0$ are defined in Definition~\ref{def:dr}.

\textbf{TMLE} uses the efficient influence function (EIF), which has the same form as the DR influence function but with targeted predictions $\hat{Q}^*$.

\subsection{Bootstrap Alternative}

Bootstrap is a valid alternative for variance estimation in weighting estimators:

\begin{enumerate}
    \item Resample $(Y_i, T_i, X_i)$ with replacement
    \item Re-estimate propensity, outcome models, and ATE
    \item Repeat $B$ times (typically $B = 1000$)
    \item Use standard deviation of bootstrap estimates as SE
\end{enumerate}

Unlike PSM with replacement, bootstrap is valid for IPW/DR/TMLE because these estimators are smooth functions of sample moments.

% ============================================================================
\section{Implementation}
\label{sec:weighting_implementation}
% ============================================================================

\subsection{IPW Implementation}

\begin{minted}{python}
def ipw_ate_observational(
    outcomes: np.ndarray,
    treatment: np.ndarray,
    covariates: np.ndarray,
    propensity: Optional[np.ndarray] = None,
    trim_at: Optional[Tuple[float, float]] = None,
    stabilize: bool = False,
    alpha: float = 0.05,
) -> IPWResult:
    """
    IPW ATE estimator for observational data.

    Parameters
    ----------
    outcomes : Observed outcomes (n,)
    treatment : Binary treatment (n,)
    covariates : Covariate matrix (n, p)
    propensity : Pre-computed P(T=1|X), or None to estimate
    trim_at : Percentiles for trimming, e.g., (0.01, 0.99)
    stabilize : Use stabilized weights

    Returns
    -------
    dict with estimate, se, ci_lower, ci_upper, diagnostics
    """
    # Step 1: Estimate propensity if not provided
    if propensity is None:
        prop_result = estimate_propensity(treatment, covariates)
        propensity = prop_result["propensity"]

    # Step 2: Trim extreme propensities
    if trim_at is not None:
        # Remove units with propensity outside [trim_at[0], trim_at[1]]
        mask = (propensity >= trim_at[0]) & (propensity <= trim_at[1])
        outcomes, treatment, propensity = outcomes[mask], treatment[mask], propensity[mask]

    # Step 3: Clip for numerical stability
    propensity = np.clip(propensity, 1e-6, 1 - 1e-6)

    # Step 4: Compute weights
    if stabilize:
        p_t = treatment.mean()
        weights = np.where(treatment == 1, p_t / propensity, (1 - p_t) / (1 - propensity))
    else:
        weights = np.where(treatment == 1, 1 / propensity, 1 / (1 - propensity))

    # Step 5: Compute IPW estimate
    ate = (treatment * outcomes * weights).sum() / (treatment * weights).sum() \
        - ((1 - treatment) * outcomes * weights).sum() / ((1 - treatment) * weights).sum()

    return {"estimate": ate, ...}
\end{minted}

\subsection{Doubly Robust Implementation}

\begin{minted}{python}
def dr_ate(
    outcomes: np.ndarray,
    treatment: np.ndarray,
    covariates: np.ndarray,
    propensity: Optional[np.ndarray] = None,
    outcome_models: Optional[Dict] = None,
    alpha: float = 0.05,
) -> DRResult:
    """
    Doubly robust ATE estimator.

    DR formula:
    ATE = (1/n) * sum[
        T/e(X) * (Y - mu1(X)) + mu1(X)
        - (1-T)/(1-e(X)) * (Y - mu0(X)) - mu0(X)
    ]
    """
    # Step 1: Estimate nuisance functions
    if propensity is None:
        propensity = estimate_propensity(treatment, covariates)["propensity"]

    if outcome_models is None:
        outcome_result = fit_outcome_models(outcomes, treatment, covariates)
        mu0, mu1 = outcome_result["mu0_predictions"], outcome_result["mu1_predictions"]

    # Step 2: Clip propensity
    propensity = np.clip(propensity, 1e-6, 1 - 1e-6)

    # Step 3: Compute DR components
    treated_term = treatment / propensity * (outcomes - mu1) + mu1
    control_term = (1 - treatment) / (1 - propensity) * (outcomes - mu0) + mu0

    # Step 4: Estimate ATE
    ate = np.mean(treated_term - control_term)

    # Step 5: Influence function variance
    influence = treated_term - control_term - ate
    se = np.sqrt(np.var(influence) / len(outcomes))

    return {"estimate": ate, "se": se, ...}
\end{minted}

\subsection{TMLE Implementation}

\begin{minted}{python}
def tmle_ate(
    outcomes: np.ndarray,
    treatment: np.ndarray,
    covariates: np.ndarray,
    max_iter: int = 100,
    tol: float = 1e-6,
) -> TMLEResult:
    """
    Targeted Maximum Likelihood Estimation for ATE.
    """
    # Step 1: Initial estimation
    propensity = estimate_propensity(treatment, covariates)["propensity"]
    outcome_result = fit_outcome_models(outcomes, treatment, covariates)
    Q0, Q1 = outcome_result["mu0_predictions"], outcome_result["mu1_predictions"]
    Q_obs = np.where(treatment == 1, Q1, Q0)

    # Step 2: Clever covariate
    H = treatment / propensity - (1 - treatment) / (1 - propensity)

    # Step 3: Targeting loop
    Q_star = Q_obs.copy()
    epsilon_total = 0.0

    for _ in range(max_iter):
        # Check convergence: score equation satisfied?
        score = np.mean(H * (outcomes - Q_star))
        if abs(score) < tol:
            break

        # Fit fluctuation: Y ~ epsilon * H + offset(Q_star)
        epsilon = np.sum(H * (outcomes - Q_star)) / np.sum(H ** 2)
        Q_star = Q_star + epsilon * H
        epsilon_total += epsilon

    # Step 4: Apply total fluctuation to both treatment levels
    H1, H0 = 1 / propensity, -1 / (1 - propensity)
    Q1_star = Q1 + epsilon_total * H1
    Q0_star = Q0 + epsilon_total * H0

    # Step 5: Compute ATE
    ate = np.mean(Q1_star - Q0_star)

    return {"estimate": ate, "converged": True, ...}
\end{minted}

% ============================================================================
\section{Monte Carlo Validation}
\label{sec:weighting_validation}
% ============================================================================

We validate all three estimators using our standard Monte Carlo protocol: 5,000 replications, $n = 200$, true ATE $= 2.0$.

\subsection{Data Generating Process}

\begin{minted}{python}
def generate_confounded_dgp(n: int, true_ate: float = 2.0):
    """Generate observational data with confounding."""
    X = np.random.normal(0, 1, (n, 2))

    # Propensity: depends on X
    logit = 0.8 * X[:, 0] + 0.5 * X[:, 1]
    e_X = 1 / (1 + np.exp(-logit))
    T = np.random.binomial(1, e_X)

    # Outcome: depends on T and X (confounding)
    Y = true_ate * T + 0.5 * X[:, 0] + 0.3 * X[:, 1] + np.random.normal(0, 1, n)

    return Y, T, X
\end{minted}

\subsection{Results}

\begin{center}
\begin{tabular}{lcccc}
\toprule
Estimator & Bias & RMSE & Coverage & Mean SE \\
\midrule
IPW & 0.04 & 0.32 & 94.8\% & 0.31 \\
IPW (stabilized) & 0.03 & 0.28 & 95.2\% & 0.27 \\
Doubly Robust & 0.02 & 0.24 & 95.6\% & 0.23 \\
TMLE & 0.01 & 0.23 & 95.4\% & 0.22 \\
\bottomrule
\end{tabular}
\end{center}

\textbf{Key findings}:
\begin{enumerate}
    \item All estimators achieve target coverage (93--97\%)
    \item TMLE and DR have lowest bias and RMSE
    \item Stabilized IPW improves on standard IPW
    \item Mean SE tracks RMSE well (variance estimation works)
\end{enumerate}

\subsection{Robustness to Misspecification}

We also test robustness when one model is wrong:

\begin{center}
\begin{tabular}{lccc}
\toprule
Scenario & IPW Bias & DR Bias & TMLE Bias \\
\midrule
Both correct & 0.04 & 0.02 & 0.01 \\
Propensity correct, outcome wrong & 0.04 & 0.03 & 0.02 \\
Propensity wrong, outcome correct & 0.35 & 0.03 & 0.02 \\
Both wrong & 0.35 & 0.18 & 0.15 \\
\bottomrule
\end{tabular}
\end{center}

This confirms double robustness: DR and TMLE remain nearly unbiased when only one model is correct.

% ============================================================================
\section{Practical Guidance}
\label{sec:weighting_guidance}
% ============================================================================

\subsection{Method Selection}

\medskip
\noindent\textbf{Decision Tree}:
\begin{enumerate}
    \item \textbf{Good overlap, simple models}: IPW often sufficient
    \item \textbf{Uncertain about model specification}: Use DR or TMLE
    \item \textbf{Need efficiency}: Use TMLE
    \item \textbf{Bounded outcomes}: TMLE can enforce bounds
    \item \textbf{Extreme propensities}: Trim first, then any method
\end{enumerate}

\subsection{Common Pitfalls}

\begin{enumerate}
    \item \textbf{Ignoring positivity violations}: Check propensity distribution before estimating
    \item \textbf{Not trimming extreme weights}: Single observations can dominate IPW
    \item \textbf{Overfitting nuisance models}: Use cross-fitting for ML-based estimators
    \item \textbf{Wrong variance estimator}: Use influence function or bootstrap, not naive formula
\end{enumerate}

\subsection{Diagnostics Checklist}

Before reporting weighting estimates:
\begin{enumerate}
    \item Check propensity score overlap (histogram by treatment)
    \item Examine weight distribution (max, effective $n$)
    \item Verify covariate balance after weighting
    \item Compare IPW, DR, TMLE estimates (should be similar if models reasonable)
    \item Run sensitivity analysis (Chapter~\ref{ch:sensitivity})
\end{enumerate}

% ============================================================================
\section{Summary}
\label{sec:weighting_summary}
% ============================================================================

\subsection{Key Results}

\begin{enumerate}
    \item \textbf{IPW}: Consistent under correct propensity model; simple but sensitive to extreme weights

    \item \textbf{Doubly Robust}: Consistent if \emph{either} propensity \emph{or} outcome model correct; recommended default

    \item \textbf{TMLE}: Achieves efficiency bound with better finite-sample properties; iterative targeting

    \item \textbf{Variance}: Use influence function approach; bootstrap valid for all three
\end{enumerate}

\subsection{Notation Reference}

\begin{center}
\begin{tabular}{ll}
\toprule
Symbol & Meaning \\
\midrule
$e(X)$ & Propensity score $\Prob(T=1 \mid X)$ \\
$\mu_t(X)$ & Outcome model $\E[Y \mid T=t, X]$ \\
$H$ & Clever covariate (TMLE) \\
$\hat{Q}^*$ & Targeted outcome predictions \\
IF & Influence function \\
EIF & Efficient influence function \\
\bottomrule
\end{tabular}
\end{center}

\subsection{Key References}

\begin{itemize}
    \item \cite{robins1994estimation}: Original doubly robust estimation
    \item \cite{bang2005doubly}: DR in causal inference
    \item \cite{hirano2003efficient}: Efficient IPW
    \item \cite{vdl2011tmle}: TMLE methodology
\end{itemize}
